\documentclass[11pt]{article}
\usepackage[margin=0.85in]{geometry}
\usepackage{amsmath,booktabs,xurl,hyperref}
\usepackage[T1]{fontenc}
\setlength{\emergencystretch}{2em}
\title{One Affine Prediction and Two Lexicographic Solves}
\author{Pan Dynamics Collision Avoidance Research}
\date{October 1, 2026}
\begin{document}
\maketitle

\section{Implemented behavior}
At the explicit project direction, the online controller now builds one affine
model and solves exactly two problems on each successful call:
\[
 J_B^*=\min\sum_{i=0}^{N-1}\xi_i,\qquad
 \min\rho\quad\text{subject to}\quad\sum_{i=0}^{N-1}\xi_i
 \le J_B^*+\varepsilon_{\mathrm{lex}}.
\]
The shared constraints include affine dynamics, actuator magnitude and slew,
state bounds, affine rectangle separation, a first-step relaxed quadratic CLF
and the existing augmented free-pose terminal core. The CLF error reference
remains the given path and desired speed. The configured 6-mm buffer is
retained in the affine rows; road constraints remain absent.

The quadratic CLF and terminal ellipsoid are represented by second-order
cones, so both stages use \texttt{coneprog}, rather than a linearly constrained
QP. There is no outer nonlinear iteration. The second stage runs even when
PCBF slack is zero. No tracking/input cost can trade against either slack,
and there is no third cost tie-break. Nonunique minimizing trajectories remain
possible.

The previous affine trajectory supplies the next linearization when usable;
otherwise a single moving-target flow rollout supplies it. Affine dynamics
retain the reference defect, so shifted affine states need not satisfy the
nonlinear map. The horizon is fixed at
$\min(M_{\max},N+\lceil T_{\mathrm{completion}}/h\rceil)$, without admission-driven
extension. The existing completion allowance is retained, giving 68 holds at
$N=8$, $T_{\mathrm{completion}}=3$ s and $h=0.05$ s.

Removed online work includes nonlinear candidate replay/admission, sequential
convexification, restoration solves, zero-CLF trial solves, nonlinear correction
QPs, endpoint polishing, predictive-cost refinement and retained-control
fallback. Only the second-stage first input is issued. A failed stage or
exhausted budget reports \texttt{noOptimizationSolution}. The previous plan is
not an executable witness. Controller state version advances to 61.

\section{Scope of the result}
Returned states and collision margins are affine-model quantities. Positive
optimized PCBF slack is explicitly allowed and reported; it is not a
collision-free claim. The second-stage numerical cap can also permit a tiny
positive slack when the primary optimum is zero. Metadata distinguishes this
scope from nonlinear validation and nonlinear-plant recursive feasibility.
Zero CLF slack establishes the imposed affine one-step inequality, not a new
proof of nonlinear convergence.

The existing terminal core construction and anchor-based separation/departure
rows remain. Encounter exit is selected once on the anchor. These calculations
are not post-solve nonlinear candidate checks. Their geometry and model tests
are retained, without transferring their isolated mathematical results into a
claim about realized nonlinear affine-prediction endpoints.

The scenario driver exports both optimization stages and raw affine residuals.
The offline Python audit accepts those numerical residuals only under the
declared affine tolerance, while physical sampled clearance remains a separate
strict check. Historical nonlinear-admission exports retain their original
interpretation. The retired option to run later frames without optimization is
removed; old continuation state cannot resume the new formulation.

\section{Validation actually performed}
\begin{center}
\begin{tabular}{lrrr}\toprule
Check & Passed & Failed & Incomplete \\\midrule
Two-stage and affine CLF tests & 13 & 0 & 0 \\
Model, geometry, terminal, interface and config tests & 242 & 0 & 0 \\
Python export/audit tests & 18 & 0 & 0 \\\bottomrule
\end{tabular}
\end{center}
The 13 targeted tests check both objectives, the primary cap, affine dynamics
including shifted-reference defects, exact first-input output, finite slew and
magnitude bounds, positive optimized slack, deadline failure without fallback,
and zero-slack affine CLF dissipation on straight/circular references at
8 and 15 m/s. The 242 related checks include eight existing compiled-kernel
equivalence/domain tests (random streams 7 and 11). Existing tests of retired
execution-witness behavior and unconditional nonlinear dissipation were replaced
by the explicit affine-contract tests; no scenario success is inferred from
that replacement.

MATLAB Code Analyzer reports no syntax errors. It retains sparse-assembly and
array-growth advisories in the solver and an obsolete suppression in the
pre-existing geometry implementation. Exact findings and test results are in
\path{report/TWO_STAGE_CONTROLLER_20261001/}. The two-stage tests and regression
checks ran with \texttt{matlab -singleCompThread -batch}; no campaign runner was
invoked. A preliminary single no-target call also completed both stages, but
its cold-call timing is not a runtime benchmark.

The earlier fourteen-case campaign remains stopped at the user's request.
This change establishes neither a new avoidance success rate nor a 100-ms
runtime bound. The configurable soft shared budget remains 5 s; a cap of two
optimizer calls is not a measured execution deadline.

\section{Reproducibility and preserved work}
The base project commit is
\path{aa21638f136804f4c7573bc87b81f9f26c37ba2d}. The implementation preserves and
includes the pre-existing in-scope 50-ms mesh/interpolated-midpoint and optional
model/geometry-kernel source changes, build script and equivalence tests. Those
changes were already present before this simplification and are not presented
as newly developed here. Compiled binaries and the nested \path{solver/} tree
are excluded from the project commit. Unrelated estimator, manuscript,
perception and report changes are preserved outside this commit.

The artifact bundle contains compact test tables, analyzer findings, actual
commands, native dependency hashes and source hashes. Existing optional kernels
were used and checked; their build script was not rerun. Full campaign,
closed-loop recovery, maximum-runtime and observer/perception regressions were
not run. The current architecture is documented in
\path{controller/PCBF_CLF_ARCHITECTURE.md}.
\end{document}
